\section{A finite determinant toolkit}
\label{sec:determinants}

Write $[n]=\{0,\ldots,n-1\}$. A permutation is a list containing each
element of $[n]$ once, and its sign is $(-1)^{\operatorname{inv}(\sigma)}$,
where an inversion is a pair $i<j$ with $\sigma(i)>\sigma(j)$. Define
\[
 \det M=\sum_{\sigma\in S_n}\operatorname{sgn}(\sigma)
                              \prod_{i<n}m_{i,\sigma(i)}.
 \tag{D}
\]
In arithmetic the sum can run through $n^n$ candidate lists and assign
zero to every list that is not a permutation. The determinant of the
empty matrix is $1$.

\begin{lemma}[Permutation signs]
Swapping two positions in a permutation reverses its sign. Moreover,
$\operatorname{sgn}(\sigma\tau)=\operatorname{sgn}(\sigma)
\operatorname{sgn}(\tau)$ and
$\operatorname{sgn}(\sigma^{-1})=\operatorname{sgn}(\sigma)$.
\end{lemma}
\begin{proof}
An adjacent interchange changes the inversion count by exactly $1$ in
absolute value. Swapping positions $i<j$ is a product of
$2(j-i)-1$ adjacent interchanges, hence changes parity. Bubble sorting a
list uses at most $n(n-1)/2$ adjacent interchanges and reduces its
inversion count to zero. Applying the same adjacent interchanges to
the positions in any other permutation multiplies its sign by the same
factor. This proves the product formula, then the inverse formula by
applying it to $\sigma\sigma^{-1}$. Every procedure and parity invariant
here is bounded in $n$.
\end{proof}

\begin{lemma}[Multilinearity, alternation, and transpose]
The determinant is additive and scalar-linear in each individual row or
column. It vanishes when two rows, or two columns, are equal.
Also $\det(M^{\mathsf T})=\det M$ and $\det I_n=1$.
\end{lemma}
\begin{proof}
Each summand of (D) contains one entry from a prescribed row and one
from a prescribed column, so distributivity proves multilinearity.
If rows $i$ and $j$ agree, pair $\sigma$ with the permutation obtained by
swapping its values at $i,j$. The products agree and the signs are
opposite. There are no fixed points. Sum once over the smaller member
of each pair to get zero. This argument never infers $z=0$ from $2z=0$.
The column proof is the same, or follows by transposition. Reindexing
by $\sigma^{-1}$ proves transposition invariance. For the identity matrix
only the identity permutation contributes, and its contribution is $1$.
\end{proof}

For $n\ge1$, write $M^{\widehat i,\widehat j}$ for the matrix obtained by
deleting row $i$ and column $j$, retaining their original orders. Define
\[
 C_{ij}=(-1)^{i+j}\det M^{\widehat i,\widehat j},\qquad
 (\operatorname{adj}M)_{ji}=C_{ij}.
\]

\begin{lemma}[Cofactor expansion and the two adjugate identities]
\label{lem:adjugate}
For every row $i$ and column $j$,
\[
 \det M=\sum_{k<n}m_{ik}C_{ik}
       =\sum_{k<n}C_{kj}m_{kj},
 \qquad
 M\operatorname{adj}M=\operatorname{adj}M\,M=(\det M)I_n.
\]
\end{lemma}
\begin{proof}
Partition the permutations according to the value $\sigma(i)=k$.
Delete the position $i$ and the value $k$, then relabel the remaining
ordered sets by $[n-1]$. This is a bijection with $S_{n-1}$; its inverse
inserts the missing position and value. Moving row $i$ to the first
position uses $i$ adjacent interchanges, and moving column $k$ to the
first position uses $k$; the surviving sign is therefore
$(-1)^{i+k}$ times the smaller permutation sign. The corresponding
products agree after extracting $m_{ik}$. This proves row expansion;
transpose gives column expansion.

The $(i,j)$ entry of $M\operatorname{adj}M$ is
$\sum_k m_{ik}C_{jk}$. It is the expansion along row $j$ of the matrix
obtained by replacing row $j$ of $M$ with row $i$. If $i=j$, it is
$\det M$; otherwise the replacement matrix has two equal rows and its
determinant is zero. Column replacement proves the other identity.
For $n=1$ the deleted determinant is $1$, so the proof includes that case.
\end{proof}

\begin{lemma}[Multiplicativity]
\label{lem:detmul}
For square matrices over a commutative ring,
$\det(AB)=\det A\det B$.
\end{lemma}
\begin{proof}
Column $j$ of $AB$ is $\sum_i b_{ij}A_i$, where $A_i$ is column $i$ of
$A$. Repeated multilinearity expands its determinant as a sum indexed
by maps $f:[n]\to[n]$, with term
\[
                 \left(\prod_j b_{f(j),j}\right)
                        \det(A_{f(0)},\ldots,A_{f(n-1)}).
\]
A noninjective $f$ repeats a column and contributes zero. Since its
domain and codomain have the same finite size, an injective $f$ is a
permutation: the occupancy counts $\nu_i=|f^{-1}(i)|$ all belong to
$\{0,1\}$ and sum to $n$, so a missing value would make their sum at
most $n-1$. Reordering its columns gives
$\operatorname{sgn}(f)\det A$. Factoring out $\det A$ leaves the column
form of (D) for $B$, proving the formula. The expansion has $n^n$ terms,
and is valid in the zero ring as well.
\end{proof}

\begin{corollary}[Change of scalars and similarity]
Every ring homomorphism preserves determinants and characteristic
polynomials. If $P,Q$ are supplied with $PQ=QP=I$, then
$\chi_{PAQ}=\chi_A$ and $f(PAQ)=P f(A)Q$ for every scalar polynomial $f$.
\end{corollary}
\begin{proof}
A homomorphism preserves the finite sum and products in (D).
The equality $tI-PAQ=P(tI-A)Q$ holds in the polynomial matrix ring, and
$\det P\det Q=1$ by Lemma~\ref{lem:detmul}. This proves the polynomial
identity. Induction gives $(PAQ)^k=PA^kQ$ for all $k$, including $k=0$;
summing proves the evaluation identity.
\end{proof}

\section{Cayley--Hamilton without division or spectral theory}
\label{sec:main}

Let $n\ge1$ and $A\in M_n(R)$. Work with finite polynomial lists in a
central indeterminate $t$. Set
\[
 \chi_A(t)=\det(tI-A)=\sum_{j=0}^n c_jt^j,
 \qquad B(t)=\operatorname{adj}(tI-A)=\sum_{j=0}^{n-1}B_jt^j.
\]
The degree bound for $B$ follows because every cofactor is a product
sum with $n-1$ linear factors. In $\chi_A$ a term of degree $n$ must
select $t$ from all $n$ rows. Only the identity permutation permits this,
and the coefficient is $1$. Thus $c_n=1$. At $t=0$,
$c_0=\det(-A)=(-1)^n\det A$.

\begin{theorem}[Coefficient telescoping]
\label{thm:CH}
For every square matrix over every commutative ring,
\[
       \boxed{\;A^n+c_{n-1}A^{n-1}+\cdots+c_1A+c_0I=0.\;}
\]
\end{theorem}
\begin{proof}
Lemma~\ref{lem:adjugate}, applied over the polynomial ring, gives
\[
                    (tI-A)B(t)=\chi_A(t)I.
\]
Comparing the finite coefficient lists gives
\begin{align*}
 B_{n-1}&=I,\\
 B_{j-1}-AB_j&=c_jI &&(1\le j\le n-1),\\
 -AB_0&=c_0I.
\end{align*}
Multiply the equation indexed by $j$ on the left by $A^j$ and add,
including the leading equation multiplied by $A^n$. The left side is
\[
 -AB_0+\sum_{j=1}^{n-1}
       \bigl(A^jB_{j-1}-A^{j+1}B_j\bigr)+A^nB_{n-1}=0;
\]
every summand cancels with its neighbor. The right side is precisely
$\sum_{j=0}^n c_jA^j$. In detail, the partial-sum invariant is
$\sum_{j=0}^k c_jA^j=-A^{k+1}B_k$ for $k<n$, followed by
$B_{n-1}=I$. This is the quantifier-free invariant identified in
Section~\ref{sec:induction}; it proves Theorem~\ref{thm:arithmetic} after
generic polynomial coding. No assertion about eigenvalues is used.

For $n=0$, the matrix algebra has one element, its zero and identity
coincide, and $\chi_A(t)=1$. Its evaluation is the empty identity matrix,
which is also zero. Thus the theorem includes dimension zero.
\end{proof}

\begin{theorem}[Adjugate as a polynomial]
\label{thm:adjpoly}
For $n\ge1$ and $0\le j<n$,
\[
 B_j=\sum_{k=j+1}^n c_k A^{k-1-j}.
\]
In particular,
\[
 \operatorname{adj}(-A)=A^{n-1}+c_{n-1}A^{n-2}+\cdots+c_1I,
 \qquad
 \operatorname{adj}A=(-1)^{n-1}
              \sum_{k=1}^{n}c_kA^{k-1}.
\]
\end{theorem}
\begin{proof}
Start with $B_{n-1}=c_nI=I$. If the formula holds for $B_j$, the
coefficient recurrence gives $B_{j-1}=AB_j+c_jI$, which is exactly the
next formula. This is a descending bounded recursion, or an ordinary
induction on $n-1-j$. The constant coefficient of the adjugate is
$\operatorname{adj}(-A)$, and each of its minors has $n-1$ negated
factors, giving the second identity. In particular every $B_j$ commutes
with $A$; this is a proved conclusion, not an assumption used above.
\end{proof}

\subsection{A second proof and the evaluation pitfall}

For a polynomial $F(t)=\sum F_it^i$ with coefficients in an associative
ring, define right evaluation by $F(A)=\sum F_iA^i$. It always preserves
addition, but generally does not preserve multiplication. If every
coefficient of $G$ commutes with $A$, however,
\[
 F(A)G(A)=\sum_{i,j}F_iA^iG_jA^j
         =\sum_{i,j}F_iG_jA^{i+j}=(FG)(A).
\]
This proves the required restricted substitution rule by a finite
calculation. In particular, evaluation is a ring homomorphism on
polynomials whose coefficients commute with $A$, even if those
coefficients do not commute with one another.

The two adjugate identities imply
$(tI-A)B(t)=B(t)(tI-A)$, hence $AB_j=B_jA$ by coefficient comparison.
The restricted substitution rule now legitimately evaluates
$(tI-A)B(t)=\chi_A(t)I$ at $A$, giving $0=\chi_A(A)$. This is a second,
fully finite proof. The stronger formula of Theorem~\ref{thm:adjpoly}
places these coefficients in the commutative algebra generated by $A$.

There is a concrete obstruction to unrestricted substitution. Take
\[
 A=\begin{pmatrix}0&1\\0&0\end{pmatrix},\qquad
 D=\begin{pmatrix}0&0\\1&0\end{pmatrix},\qquad F(t)=tI,\quad G(t)=D.
\]
Then $(FG)(A)=DA$, whereas $F(A)G(A)=AD$; the two products are different
diagonal idempotents. Thus a proof that treats evaluation on all
matrix-coefficient polynomials as a homomorphism is invalid.
Nor can one replace the scalar variable in $\det(tI-A)$ by a matrix
inside the entries and call the resulting expression $\det(A-A)$:
the types and the surrounding matrix sizes have changed.

An additional arithmetic diagnostic is provided by the permanent. For
$H=\left(\begin{smallmatrix}1&2\\3&4\end{smallmatrix}\right)$,
\[
 \operatorname{perm}(tI-H)=t^2-5t+10,
 \qquad H^2-5H+10I=12I\ne0.
\]
The same putative substitution proof would incorrectly annihilate this
permanent polynomial. Our proof specifically uses the sign cancellation
in the determinant.

\subsection{Inverse matrices, including the exact hypothesis}

\begin{theorem}[Explicit inverse from a unit determinant]
\label{thm:inverse}
For $n\ge1$, if $u\det A=1$ is supplied, then
\[
 A^{-1}=u\operatorname{adj}A
       =u(-1)^{n-1}
                (A^{n-1}+c_{n-1}A^{n-2}+\cdots+c_1I)
\]
is a two-sided inverse. Conversely a supplied inverse $V$ gives the
unit witness $\det V$ for $\det A$.
\end{theorem}
\begin{proof}
Multiply each adjugate identity by the scalar $u$ for the first claim.
The explicit polynomial expression follows from
Theorem~\ref{thm:adjpoly}. Conversely $AV=I$ gives
$\det A\det V=\det I=1$ by determinant multiplicativity. All witnesses
are explicit finite expressions.
\end{proof}
Over a field a nonzero determinant has an inverse, so this recovers the
usual statement. Over $\mathbb Z$, the matrix $(2)$ has nonzero
determinant and no inverse over $\mathbb Z$; nonzeroness is not the
correct ring hypothesis. The theorem remains valid in the zero ring,
where the required equations have their usual algebraic meanings.

\subsection{The displayed small examples}

For $A=(a)$ one has $\chi_A(t)=t-a$ and $A-aI=0$. For
$A=\left(\begin{smallmatrix}a&b\\c&d\end{smallmatrix}\right)$,
\[
 \chi_A(t)=t^2-(a+d)t+(ad-bc),\qquad
 A^2=\begin{pmatrix}a^2+bc&ab+bd\\ca+dc&cb+d^2\end{pmatrix}.
\]
Subtracting $(a+d)A$ and adding $(ad-bc)I$ cancels each of the four
entries, using only commutativity and distributivity.
For the numerical matrix $H$ above this gives
\[
 H^2=5H+2I,\qquad H^3=27H+10I,\qquad H^4=145H+54I.
\]
Each successive equality is obtained by multiplying the previous one
by $H$ and replacing $H^2$ by $5H+2I$. Over $\mathbb Q$,
\[
 H^{-1}=\tfrac12(H-5I),\qquad
 H^{-2}=\tfrac14(-5H+27I).
\]
The first identity follows from $H(H-5I)=2I$; squaring it and reducing
$H^2$ proves the second. Thus all of the article's elementary numerical
examples have exact integer or rational proofs.

